\setcounter{section}{4}
\Section{Применение производной к различным вопросам математического анализа}
\setcounter{subsection}{4}
\Subsection{Выпуклости и точки изгиба}

\begin{definition}
Функция \(f(x)\) называется \textit{выпуклой вверх (вниз) на промежутке~\(I\)},
если
\[
    \forall x_1, x_2 \in I, x_1 \neq x_2 \quad f\left(\frac{x_1 + x_2}{2}
        \right) > \frac{f(x_1) + f(x_2)}{2} \quad \left(f\left(\frac{x_1 + x_2}
        {2}\right) < \frac{f(x_1) + f(x_2)}{2}\right)
\]
\end{definition}

В этих случаях можно говорить о \textit{строгой выпуклости}. Если применять
нестрогие неравенства~--- \textit{нестрогая выпуклость}.

\begin{definition}
Точка \(x_0\) называется \textit{точкой перегиба}~\(f(x)\), если \(\exists
f'(x)\) конечная, \(+\infty\) или \(-\infty\) и на промежутках \((x_0 - \delta,
x_0]\) и \([x_0, x_0 + \delta)\), \(\delta > 0\) выпуклость в разные стороны.
\end{definition}

\begin{statement}
Если существует \(f''(x_0)\), то
\[
    f''(x_0) = \lim_{t \rightarrow 0} \frac{f(x_0 + t) + f(x_0 - t) - 2f(x_0)}
        {t^2}.
\]
\end{statement}

\begin{proof}
\begin{gather}
    f(x_0 + t) = f(x_0) + f'(x_0) \cdot t + \frac{f''(x_0)}{2} \cdot t^2 +
        o(t^2), \; t \rightarrow 0\\
    f(x_0 - t) = f(x_0) - f'(x_0) \cdot t + \frac{f''(x_0)}{2} \cdot t^2 +
        o(t^2)
\end{gather}
сложим два этих равенства:
\[
    f(x_0 + t) + f(x_0 - t) = 2f(x_0) + f''(x_0) \cdot t^2 + o(t^2)
\]
отсюда следует:
\[
    \frac{f(x_0 + t) + f(x_0 - t) - 2f(x_0)}{t^2} = f''(x_0) + o(1)
\]
\end{proof}

\begin{note}
Такой предел может существовать, а \(f''(x_0)\)~--- нет. Например, \(f(x)\)~---
нечетная функция, \(x_0 = 0\):
\[
    \frac{f(t) + f(-t) - 2 \cdot 0}{t^2} \equiv 0
\]
\end{note}

\begin{theorem}[Необходимое условие выпуклости]
Пусть функция \(f(x)\) имеет конечную~\(f''(x)\) на интервале~\(I\), конечном
или бесконечном. Тогда, если \(f(x)\) выпукла вверх (вниз) на \(I\), то \(f''(x)
< 0 \; (f''(x) > 0)\)
\end{theorem}

\begin{proof}
Пусть \(x_0 \in I, \; x_1 = x_0 + t, \; x_2 = x_0 - t, \; x_1, x_2 \in I, \; t
\neq 0\). Доказательство проведём для выпуклости вверх:
\begin{gather}
    f\left(\frac{x_1 + x_2}{2}\right) > \frac{f(x_1) + f(x_2)}{2}\\
    f(x_0) > \frac{f(x_0 + t) + f(x_0 - t)}{2}\\
    f(x_0 + t) + f(x_0 - t) - 2f(x_0) = f(x_1) + f(x_2) -
        2f\left(\frac{x_1 + x_2}{2}\right) < 0 \; (\text{при достаточно малых}
        \; t)\\
    \frac{f(x_0 + t) + f(x_0 - t) - 2f(x_0)}{t^2} < 0\\
    \lim_{t \rightarrow 0}{\frac{f(x_0 + t) + f(x_0 - t) - 2f(x_0)}{t^2}} \leq 0
        \Rightarrow f''(x_0) \leq 0
\end{gather}
\end{proof}

\begin{note}
\(f''(x) \leq 0 \; (\geq 0)\) независимо от того, строгую или нестрогую
выпуклость мы рассматриваем.
\end{note}

\begin{example}
\(f(x) = x^4\), строго выпукла вниз, \(f''(0) = 0\), \(f''(x) = 12x^2 \geq 0\).
\end{example}

\begin{theorem}[Достаточное условие выпуклости]
Пусть \(f(x)\) непрерывна на промежутке~\(I\) и существует конечная~\(f''(x)\)
во внутренних точках~\(I\). Тогда, если \(f''(x) > 0\) во внутренних
точках~\(I\) (\(< 0\)), то \(f(x)\) выпукла вниз (вверх) на~\(I\).
\end{theorem}

\begin{proof}
Пусть \(\ds x_1, x_2 \in I, \; x_1 \neq x_2, \; x_0 = \frac{x_1 + x_2}{2}, \;
t = \frac{x_1-x_2}{2}\) \((x_1 > x_2)\). Тогда \(x_1 = x_0 + t\), \(x_2 = x_0 -
t\). Получим:
\begin{align}
    \frac{f(x_1) + f(x_2)}{2} - f\left(\frac{x_1 + x_2}{2}\right) &=
        \frac{f(x_0 + t) + f(x_0 - t) - 2f(x_0)}{2} =\\
        &= \frac{\left(f(x_0 + t) - f(x_0)\right) - \left(f(x_0) - f(x_0 - t)
        \right)}{2}
\end{align}
применив теорему Лагранжа к последнему выражению, получим
\begin{gather}
    \frac{\left(f(x_0 + t) - f(x_0)\right) - \left(f(x_0) - f(x_0 - t)\right)}
        {2} = \frac{f'(\xi_1) \cdot t - f'(\xi_2) \cdot t}{2} = \left(f'(\xi_1)
        \cdot t - f'(\xi_2)\right) \cdot \frac{t}{2},
\end{gather}
где \(\xi_1 \in (x_0, x_0 + t)\), \(\xi_2 \in (x_0 - t, x_0)\). К последней
разности применим теорему Лагранжа еще раз:
\begin{gather}
    \left(f'(\xi_1) \cdot t - f'(\xi_2)\right) \cdot \frac{t}{2} =
        f''(\xi_3) \cdot (\xi_1-\xi_2) \cdot \frac{t}{2} >0, \; \text{если} \;
        f''(x) > 0
\end{gather}
получаем:
\begin{gather}
    f\left(\frac{x_1 + x_2}{2}\right) < \frac{f(x_1) + f(x_2)}{2}
\end{gather}
Это и означает выпуклость вниз на~\(I\).
\end{proof}

\begin{note}
Если \(f'(x) > 0\), \(f''(x) > 0\), то выпуклость \textit{строгая}.
\end{note}

\begin{theorem}[Достаточное условие точки перегиба]
Пусть \(\exists f'(x_0)\), конечная, \(+ \infty\) или \(- \infty\), \(f(x)\)
непрерывна в точке \(x_0\) и \(\exists f''(x)\) в \(\overset{\circ}{U}_{\delta}
(x_0)\). Тогда, если \(f''(x) > 0\) на \((x_0 - \delta, x_0)\), \(f''(x) < 0\)
на \((x_0, x_0 + \delta)\) (или наоборот), то \(x_0\)~--- точка перегиба \(
f(x)\).
\end{theorem}

\begin{example}
\(f(x) = (x + 1)^3e^{-3|x|}\).
\begin{align}
    x > 0&: \quad f''(x) = 3(x + 1)(3x^2 - 1)e^{-3x}\\
    x < 0&: \quad f''(x) = 3(x + 1)(3x^2 + 12x + 11)e^{3x}\\
\end{align}

\textbf{\textit{[тут должен быть график пу пу пу]}}

\(x_0 = 0\) не является перегибом, хотя выпуклость по разные стороны от точки
разная (\(f'(0)\) не существует).

\(x_0 = 1/\sqrt{3}, -2 \pm 1/\sqrt{3}, -1\)~--- перегиб.

\(x_0 = -1\)~--- перегиб с горизонтальной касательной.
\end{example}

\begin{note}
Выпуклость можно ещё рассматривать как расположение по отношению к касательной.
Если \(f''(x)\) существует на \((x_0 - \delta, x_0 + \delta)\), то
\begin{gather}
    f(x) = f(x_0) + f'(x_0)(x - x_0) + \frac{f''(\xi)}{2}(x - x_0)^2
\end{gather}
\(f(x_0) + f'(x_0)(x - x_0)\) --- линейная функция соответствующая уравнению
касательной.
\begin{align}
    f''(x) > 0 \; \text{в} \; \overset{\circ}{U}_{\delta}(x_0) &\Rightarrow \;
        \text{график выше касательной,}\\
    f''(x) < 0 \; \text{в} \; \overset{\circ}{U}_{\delta}(x_0) &\Rightarrow \;
        \text{ниже касательной.}
\end{align}
\end{note}